\input{sol_pre.tex}
\input{q_2025-05-25_2.tex}
\begin{document}
\input{sol_head.tex}

\Savol{1}
Har bir tub son darajasi $a$ va $b$ orasida bölinadi; EKUK dagi daraja — kattasi, ya'ni kamida yarmi (yuqoriga yaxlitlangan): $2^2$, $3^3$, $7^1$.\\
$a=2^2\cdot3^2\cdot7=252$, $b=2\cdot3^3\cdot7=378$ da EKUK $=2^2\cdot3^3\cdot7=756$.
\Javob{A}

\Savol{2}
$a^6=144$, $b^6=216$, $c^6=64\cdot3=192$ $\Rightarrow$ $a<c<b$. Şuningdek $a>2$, $c>2$, $b<3$ $\Rightarrow$ $a+c-b>0$.
\[\frac{(b-a)+(b-c)-(c-a)-2a}{a+c-b}=\frac{-2(a+c-b)}{a+c-b}=-2.\]
\Javob{D}

\Savol{3}
Turğun suvdagi tezlik $v$: \ $\dfrac{24}{v+2}+\dfrac{24}{v-2}=9$ $\Rightarrow$ $3v^2-16v-12=0$ $\Rightarrow$ $v=6$. \ $6:2=3$.
\Javob{B}

\Savol{4}
$6x=0{,}9\cdot2y$ $\Rightarrow$ $x=0{,}3y$, ya'ni $30\%$.
\Javob{D}

\Savol{5}
$(0{,}25)^{-\frac12}=2$, \ $(0{,}5)^{-2}=4$, \ $9^{-\frac12}=\frac13$: \ $\frac23+\frac43=2$.
\Javob{A}

\Savol{6}
$\dfrac{4\left(\sqrt{13}-3\right)}{4}-\dfrac{9\left(\sqrt{13}+2\right)}9=\sqrt{13}-3-\sqrt{13}-2=-5$.
\Javob{D}

\Savol{7}
$x$ — ifoda qiymati: \ $x^3=4+3\sqrt[3]{(2+\sqrt5)(2-\sqrt5)}\cdot x=4-3x$ $\Rightarrow$ $x^3+3x-4=0$ $\Rightarrow$ $(x-1)\left(x^2+x+4\right)=0$, \ $x=1$.
\Javob{C}

\Savol{8}
$S_{14}-S_3=a_4+\dots+a_{14}=11a_9=77$ $\Rightarrow$ $a_9=7$. \ $a_4+a_{17}=2a_1+19d=17$, \ $a_1+8d=7$ $\Rightarrow$ $d=1$, $a_1=-1$.\\
$a_k=k-2=2023$ $\Rightarrow$ $k=2025$, \ $\sqrt k=45$.
\Javob{B}

\Savol{9}
$b_1\left(q^3-1\right)=52$, \ $b_1\left(q^6-1\right)=1456$ $\Rightarrow$ $q^3+1=28$, $q=3$, $b_1=2$.
\[S_6=\frac{2\left(3^6-1\right)}{3-1}=728.\]
\Javob{D}

\Savol{10}
$\dfrac1{\sqrt6+\sqrt2}-\dfrac1{\sqrt6-\sqrt2}=\dfrac{-2\sqrt2}4=-\dfrac{\sqrt2}2$; \ $\left(-\dfrac{\sqrt2}2\right)^{-2}=2$.
\Javob{A}

\Savol{11}
$\dfrac1{\sqrt k+\sqrt{k+1}}=\sqrt{k+1}-\sqrt k$ $\Rightarrow$ yiğindi $\sqrt{25}-\sqrt3=5-\sqrt3\approx3{,}27\in(3;4)$.
\Javob{C}

\Savol{12}
$\arcsin\left(-\frac12\right)=-\frac\pi6$, \ $\tg\left(-\frac\pi6\right)=-\frac{\sqrt3}3$.
\Javob{B}

\Savol{13}
$t=\sqrt x\ge0$, $t\ne2$: \ $\dfrac{3(t-2)(t-3)}{(t-2)(t+2)}>t-3$ $\Rightarrow$ $(t-3)\left(\dfrac3{t+2}-1\right)>0$ $\Rightarrow$ $\dfrac{(t-3)(1-t)}{t+2}>0$ $\Rightarrow$ $1<t<3$.\\
$1<x<9$, $x\ne4$: \ $2+3+5+6+7+8=31$.
\Javob{B}

\Savol{14}
$2^x\left(25-5^x\right)-\left(25-5^x\right)=\left(25-5^x\right)\left(2^x-1\right)\ge0$ $\Rightarrow$ $0\le x\le2$. \ $0+1+2=3$.
\Javob{C}

\Savol{15}
Şartlar: $x\ge|y|$, $y\ge0$. Birinçi tenglamani kvadratga oşiramiz: $2x+2\sqrt{x^2-y^2}=y^2$, \ $\sqrt{x^2-y^2}=6$:
\[2x+12=y^2\ \Rightarrow\ x=\frac{y^2-12}2;\qquad \frac{\left(y^2-12\right)^2}4=y^2+36\ \Rightarrow\ y^2\left(y^2-28\right)=0.\]
$y=0$ da $x=-6<0$ — mos emas; \ $y=2\sqrt7$ da $x=8\ge2\sqrt7$ — mos. Bitta juftlik.
\Javob{A}

\Savol{16}
$\lg\left(4\cdot3\cdot3^{\frac1{2x}}\right)=\lg\left(3^{\frac1x}+27\right)$, \ $t=3^{\frac1{2x}}>0$: \ $t^2-12t+27=0$ $\Rightarrow$ $t=3$ yoki $t=9$.\\
$\frac1{2x}=1$ $\Rightarrow$ $x=\frac12$; \ $\frac1{2x}=2$ $\Rightarrow$ $x=\frac14$. \ Yiğindi $\frac34$.
\Javob{B}

\Savol{17}
$\ctg x\ge0$ bölişi kerak. \ $\ctg x=0$: $x=\frac\pi2,\ \frac{3\pi}2$. \ $\sin x=\frac12$: $x=\frac\pi6$ ($\ctg>0$ — mos), $x=\frac{5\pi}6$ ($\ctg<0$ — mos emas).
\[\frac\pi6+\frac\pi2+\frac{3\pi}2=\frac{13\pi}6.\]
\Javob{C}

\Savol{18}
$2\le x\le6$. Ikkala tomon nomanfiy, 4-darajaga oşiramiz: $(x-2)^2\ge6-x$ $\Rightarrow$ $x^2-3x-2\ge0$ $\Rightarrow$ $x\ge\frac{3+\sqrt{17}}2\approx3{,}56$.\\
Butunlar: $4+5+6=15$.
\Javob{C}

\Savol{19}
$3x-x^2+4=6{,}25-(x-1{,}5)^2\le6{,}25$ $\Rightarrow$ $\sqrt{\phantom{x}}\in[0;\,2{,}5]$. Butun $a$: $0,1,2$ — $3$ ta.
\Javob{B}

\Savol{20}
$f(2x)=2x+3$, \ $2f(x-1)=2x+4$: \ $3\cdot(-1)=-3$.
\Javob{A}

\Savol{21}
$x\to\frac1x$: \ $af\left(\frac1x\right)+bf(x)=x-2024$. Birinçisini $a$ ga, ikkinçisini $b$ ga köpaytirib ayiramiz:
\[\left(a^2-b^2\right)f(x)=\frac ax-bx-2024(a-b)\ \Rightarrow\ f(x)=\frac1{a^2-b^2}\left(\frac ax-bx\right)-\frac{2024}{a+b}.\]
\Javob{B}

\Savol{22}
$c=\cos x\in\left[\frac12;\frac{\sqrt3}2\right]$: \ $y=\dfrac1{1-c}+\dfrac4c$, \ $y'_c=\dfrac1{(1-c)^2}-\dfrac4{c^2}=0$ $\Rightarrow$ $c=2(1-c)$, $c=\frac23$ (oraliqda).\\
$y\left(\frac23\right)=3+6=9$; \ çetlarda: $y\left(\frac12\right)=10$, $y\left(\frac{\sqrt3}2\right)\approx12{,}1$.
\Javob{D}

\Savol{23}
$\dfrac{\sin^6x}{\cos^8x}=\tg^6x\cdot\dfrac1{\cos^2x}$:
\[\int_0^{\frac\pi4}\tg^6x\,d(\tg x)=\frac{\tg^7x}7\Big|_0^{\frac\pi4}=\frac17.\]
\Javob{A}

\Savol{24}
Kvadrat tomoni $2r$: \ $\dfrac{8r}{2\pi r}=\dfrac4\pi$.
\Javob{A}

\Savol{25}
Oltiburçak tomoni $a$. $\angle EFG=180^\circ-120^\circ=60^\circ$. $\angle DEC=30^\circ$ ($CDE$ — teng yonli, uçidagi burçak $120^\circ$), demak $\angle CEF=120^\circ-30^\circ=90^\circ$ va $\angle FEG=90^\circ$.\\
$\triangle FEG$: $\angle G=30^\circ$ $\Rightarrow$ $FG=2EF$, $a=3$.\\
$ACEF$: $AC=CE=a\sqrt3$, $AF=EF=a$, $\angle CAF=\angle CEF=90^\circ$ — ikkita teng töğri burçakli uçburçak:
\[S=2\cdot\frac12\cdot a\cdot a\sqrt3=a^2\sqrt3=9\sqrt3.\]
\Javob{B}

\Savol{26}
$\angle C=90^\circ$ $\Rightarrow$ $AB=2R=50$. \ $r=\dfrac{a+b-c}2$ $\Rightarrow$ $a+b=60$.\\
$a^2+b^2=2500$ $\Rightarrow$ $2ab=3600-2500=1100$:
\[S=\frac{ab}2=275.\]
\Javob{D}

\Savol{27}
$\angle ANM+\angle DNC=90^\circ$ $\Rightarrow$ $\triangle MAN\sim\triangle NDC$, koeffitsiyent $\frac{MN}{NC}=\frac34$.\\
Kvadrat tomoni $s$, $DN=x$: \ $AN=\frac34DC=\frac34s$ $\Rightarrow$ $x=\frac s4$. \ $\triangle NDC$: $x^2+s^2=16$:
\[\frac{s^2}{16}+s^2=16\ \Rightarrow\ s^2=\frac{256}{17}=15\frac1{17}.\]
\Javob{A}

\Savol{28}
$AC=10$, $BD=10$, $CD=6$ $\Rightarrow$ $BC=8$. \ $S_{ABD}=\frac12\cdot AD\cdot BC=\frac12\cdot4\cdot8=16$.\\
$AE$ — $ABD$ uçburçak medianasi: \ $S_{ABE}=\frac{16}2=8$.
\Javob{D}

\Savol{29}
$MN$ — örta çiziq, şuning uçun $O$ — $BH$ ning örtasi: \ $\overrightarrow{AO}=\frac12\left(\overrightarrow{AB}+\overrightarrow{AH}\right)$.\\
Teng yonli trapetsiyada $AH=\frac{AD-BC}2=\frac{3BC-BC}2=BC$ va $\overrightarrow{AH}\upuparrows\overrightarrow{BC}$, ya'ni $\overrightarrow{AH}=\overrightarrow{BC}$:
\[\overrightarrow{AO}=\frac12\overrightarrow{AB}+\frac12\overrightarrow{BC}\ \Rightarrow\ \frac ab=1.\]
\Javob{C}

\Savol{30}
$\pi r^2+\pi rl=8\pi r^2$ $\Rightarrow$ $l=7r$, \ $h^2=48r^2=576$ $\Rightarrow$ $r^2=12$:
\[V=\frac13\pi\cdot12\cdot24=96\pi.\]
\Javob{B}

\Savol{31}
Son $a$, $b$, $c$ raqamlari bilan töliq aniqlanadi ($2025=3\cdot675$). $a\ne0$: \ $9\cdot9\cdot8=648$.
\Javob{D}

\Savol{32}
$A\cup B=[2;4]\cup[6;8]$, \ $(A\cup B)\cap C=[3;4]\cup[6;7]=D\cup E$.\\
($A\cap B=\varnothing$ bölgani uçun A da $\varnothing$, C da $C$ hosil böladi.)
\Javob{B}

\newpage
\Savol{33}
$BC\parallel AD$ ($y=2\sqrt3+2$ va $y=2$), $AD=12$, $BC=4$, balandlik $h=2\sqrt3$.\\
$B$ va $C$ ning $AD$ dagi proyeksiyalari: $AB_1=2$, $C_1D=6$; \ $AB=\sqrt{4+12}=4$, \ $CD=\sqrt{36+12}=4\sqrt3$.

\textbf{33.} $S=\dfrac{12+4}2\cdot2\sqrt3=16\sqrt3$ — \textbf{B}.

\textbf{34.} $AD$ atrofida aylantirilganda: ikkita konus (yasovçilari $4$ va $4\sqrt3$) va silindr (uzunligi $4$), radius $r=2\sqrt3$:
\[S=\pi r\left(4+4\sqrt3\right)+2\pi r\cdot4=8\sqrt3\pi+24\pi+16\sqrt3\pi=24\left(\sqrt3+1\right)\pi\ \text{— \textbf{F}}.\]

\textbf{35.} $V=\pi r^2\left(\dfrac23+4+\dfrac63\right)=12\pi\cdot\dfrac{20}3=80\pi$ — \textbf{C}.
\Javob{33 — B, \ 34 — F, \ 35 — C}

\newpage
\Savol{36}
$\left(3+2\sqrt2\right)^2=17+12\sqrt2$, \ $\left(3+2\sqrt2\right)^3=99+70\sqrt2$, \ $\left(3-2\sqrt2\right)=\left(3+2\sqrt2\right)^{-1}$.\\
$u=\left(3+2\sqrt2\right)^x>0$:
\[u^3-7u^2+8u+\frac1u=7\ \Rightarrow\ u^4-7u^3+8u^2-7u+1=0.\]
$u^2$ ga bölamiz, $s=u+\frac1u\ge2$: \ $s^2-2-7s+8=0$ $\Rightarrow$ $s^2-7s+6=0$ $\Rightarrow$ $s=6$ ($s=1<2$ — mos emas).\\
$u+\frac1u=6$ $\Rightarrow$ $u=3\pm2\sqrt2$ $\Rightarrow$ $x=\pm1$.

a) $2$ ta ildiz. \quad b) $1+(-1)=0$.
\Javob{a) 2; \quad b) 0}

\newpage
\Savol{37}
Aniqlaniş sohasi: $3x^2-7x+4\le0$ $\Rightarrow$ $(3x-4)(x-1)\le0$ $\Rightarrow$ $x\in\left[1;\frac43\right]$.\\
1) $\sqrt{7x-4-3x^2}=0$: \ $x=1$, $x=\frac43$.\\
2) $\sin x=\frac{\sqrt3}2$: \ $x=\frac\pi3\approx1{,}047\in\left[1;\frac43\right]$ ($\frac{2\pi}3$ va boşqalari oraliqdan taşqarida).

a) $3$ ta ildiz: $1$, $\frac\pi3$, $\frac43$.

b) $1\cdot\frac43=\frac43$.
\Javob{a) 3; \quad b) $\frac43$}

\newpage
\Savol{38}
$f(x)f\left(\frac1x\right)=a^2+ab\left(x^n+x^{-n}\right)+b^2$, \ $f(x)+f\left(\frac1x\right)=a\left(x^n+x^{-n}\right)+2b$.\\
Har qanday $x$ da teng: $ab=a$ $\Rightarrow$ $b=1$; \ $a^2+1=2$ $\Rightarrow$ $a=\pm1$.\\
$f(3)=a\cdot3^n+1=82$ $\Rightarrow$ $a\cdot3^n=81$ $\Rightarrow$ $a=1$, $n=4$: \ $f(x)=x^4+1$.

a) $f(2)+f(1)=17+2=19$.

b) $x^4\ge0$ $\Rightarrow$ $E(f)=[1;\infty)$.
\Javob{a) 19; \quad b) $[1;\infty)$}

\newpage
\Savol{39}
$f'(x)=\cos\frac1x\cdot\left(-\frac1{x^2}\right)=-\dfrac{\cos\frac1x}{x^2}$.

a) $f'\left(\frac1\pi\right)=-\pi^2\cos\pi=\pi^2$.

b) $f''(x)=\dfrac{2\cos\frac1x}{x^3}-\dfrac{\sin\frac1x}{x^4}$. U holda
\[g(x)=\sin\frac1x-2x\cos\frac1x+2x\cos\frac1x-\sin\frac1x=0\ \Rightarrow\ g(2025)=0.\]
\Javob{a) $\pi^2$; \quad b) 0}

\newpage
\Savol{40}
$\displaystyle\int_0^1f\,f'\,dx=\frac{f^2(1)-f^2(0)}2=0$ $\Rightarrow$ $f(0)=\pm f(1)$.\\
$\displaystyle\int_0^1f^2f'\,dx=\frac{f^3(1)-f^3(0)}3=18$ $\Rightarrow$ $f^3(1)-f^3(0)=54\ne0$, demak $f(0)=-f(1)$:
\[2f^3(1)=54\ \Rightarrow\ f(1)=3,\quad f(0)=-3.\]
a) $f(1)=3$.

b) $\displaystyle\int_0^1f^4f'\,dx=\frac{f^5(1)-f^5(0)}5=\frac{243+243}5=\frac{486}5$.
\Javob{a) 3; \quad b) $\frac{486}5$}

\newpage
\Savol{41}
$\triangle AHD$ ning asosi $AD$, balandligi $AB$: \ $S_{AHD}=\frac12S_{ABCD}$.\\
$\triangle ABE$ ning asosi $AB$, balandligi $AD$: \ $S_{ABE}=\frac12S_{ABCD}$.\\
Rasmdan: $S_{AHD}=S_1+S_2+S_{AFD}$, \ $S_{ABE}=S_{ABK}+S_2+S_{FEG}$.

a) $S_1+S_2+15=15+S_2+4$ $\Rightarrow$ $S_1=4$.

b) $S_{ABCD}=10\cdot11=110$: \ $S_2=55-15-4=36$.
\Javob{a) 4; \quad b) 36}

\newpage
\Savol{42}
$S=\frac12ah_a$ $\Rightarrow$ $a=\frac{2S}{h_a}$ va h.k.; \ $S=pr$ $\Rightarrow$ $\dfrac1r=\dfrac{a+b+c}{2S}=\dfrac1{h_a}+\dfrac1{h_b}+\dfrac1{h_c}$.

a) $\dfrac1r=\dfrac1{12}+\dfrac{13}{168}+\dfrac5{56}=\dfrac{14+13+15}{168}=\dfrac14$ $\Rightarrow$ $r=4$.

b) $a:b:c=\frac1{h_a}:\frac1{h_b}:\frac1{h_c}=14:13:15$, ya'ni $a=14k$, $b=13k$, $c=15k$.\\
Geron: $p=21k$, $S=\sqrt{21\cdot7\cdot8\cdot6}\,k^2=84k^2$. Boşqa tomondan $S=\frac12\cdot14k\cdot12=84k$ $\Rightarrow$ $k=1$:
\[S=84.\]
\Javob{a) 4; \quad b) 84}

\newpage
\Savol{43}
$6\cdot\frac{\sqrt3}4a^2=6\sqrt3$ $\Rightarrow$ $a=2$. $L$ — oltiburçak markazi, $LE=LD=2$, $LCDE$ — romb.

a) $LDE$ teng tomonli, $DF\perp BE$ va $DM$ — uning balandligi (medianasi): \ $LM=\frac{LE}2=1$.

b) $CE$ va $LD$ — $LCDE$ romb diagonallari: $CE\perp LD$ va $K$ — $LD$ ning örtasi, $LK=1$.\\
$LKNM=\triangle LKE-\triangle NME$:
\[S_{LKE}=\frac12\cdot LK\cdot LE\cdot\sin60^\circ=\frac12\cdot1\cdot2\cdot\frac{\sqrt3}2=\frac{\sqrt3}2.\]
$\triangle NME$: $\angle M=90^\circ$, $\angle E=\angle CEB=30^\circ$, $ME=1$ $\Rightarrow$ $MN=\frac1{\sqrt3}$, \ $S_{NME}=\frac12\cdot1\cdot\frac1{\sqrt3}=\frac{\sqrt3}6$.
\[S_{LKNM}=\frac{\sqrt3}2-\frac{\sqrt3}6=\frac{\sqrt3}3.\]
\Javob{a) 1; \quad b) $\frac{\sqrt3}3$}

\newpage
\Savol{44}
$\pi R^2H=\frac13\pi R^2H+\frac43\pi R^3$ $\Rightarrow$ $\frac23H=\frac43R$ $\Rightarrow$ $H=2R$.

a) $H=10$ $\Rightarrow$ $R=5$: \ $V=\pi\cdot25\cdot10=250\pi$ cm$^3=\frac\pi4$ l $\approx0{,}785$ l.

b) $4\pi R^2=144\pi$ $\Rightarrow$ $R=6$, $H=12$: \ $V=\pi\cdot36\cdot12=432\pi$.
\Javob{a) $\frac\pi4\approx0{,}785$ l; \quad b) $432\pi$}

\newpage
\Savol{45}
Sektorlar burçaklari $120^\circ$ va $240^\circ$. Konus yasovçisi $R$, asos aylanasi uzunligi — yoy uzunligi:
\[2\pi r_1=\frac{2\pi R}3\Rightarrow r_1=\frac R3,\ h_1=\frac{2\sqrt2}3R;\qquad 2\pi r_2=\frac{4\pi R}3\Rightarrow r_2=\frac{2R}3,\ h_2=\frac{\sqrt5}3R.\]
a) $\dfrac{V_1}{V_2}=\dfrac{r_1^2h_1}{r_2^2h_2}=\dfrac14\cdot\dfrac{2\sqrt2}{\sqrt5}=\dfrac{\sqrt2}{2\sqrt5}=\dfrac{\sqrt{10}}{10}$.

b) Öq kesimda içki aylana radiusi $\rho=\dfrac{S}{p}=\dfrac{rh}{r+R}$:
\[\rho_1=\frac{\frac R3\cdot\frac{2\sqrt2R}3}{\frac{4R}3}=\frac{\sqrt2}6R,\qquad \rho_2=\frac{\frac{2R}3\cdot\frac{\sqrt5R}3}{\frac{5R}3}=\frac{2\sqrt5}{15}R,\qquad \frac{\rho_1}{\rho_2}=\frac{\sqrt{10}}4.\]
\[\frac{V_{sh1}}{V_{sh2}}=\left(\frac{\sqrt{10}}4\right)^3=\frac{10\sqrt{10}}{64}=\frac{5\sqrt{10}}{32}.\]
\Javob{a) $\frac{\sqrt{10}}{10}$; \quad b) $\frac{5\sqrt{10}}{32}$}

\ShowJavoblar
\end{document}
